\chapter*{Introduction}
\addcontentsline{toc}{chapter}{Introduction}
\markboth{Introduction}{Introduction}

In this chapter we introduce the topic of protocol fuzzing, motivate the choice of WebTransport as a target, and outline the structure of the thesis.

Protocol fuzzing is a dynamic testing technique used for discovering vulnerabilities in software implementations by providing invalid, unexpected, or random data as inputs to a program. In the context of network protocols, fuzzing helps uncover bugs that static analysis or traditional testing may miss. It is not a replacement for formal verification or exhaustive testing, but rather complements them by simulating realistic attacker conditions, often revealing edge cases such as buffer overflows, crashes, or logic errors arising from malformed messages or unexpected sequences. Fuzzing operates on a feedback-driven basis, where test cases are generated, executed, and mutated based on coverage or crash detection, making it particularly effective for complex systems such as network stacks. However, fuzzing is not infallible; it may miss subtle timing-based issues or require significant computational resources, and its effectiveness depends on the quality of initial seed inputs and mutation strategies.

WebTransport is a modern web API that enables low-latency, bidirectional, and multiplexed communication between web clients and servers, built on top of HTTP/3 and the QUIC transport protocol~\cite{webtransport_http3, webtransport_api}. It provides features such as streams, datagrams, and sessions for applications like gaming, live streaming, and real-time collaboration. At a high level, WebTransport establishes a session over QUIC, which handles congestion control, multiplexing, and loss recovery, while HTTP/3 provides framing for initial handshakes and capsules. The protocol manages states including connection establishment, stream creation, data transfer, and graceful shutdown, with error handling for invalid frames or capsules. Although WebTransport is still being actively developed according to the latest IETF drafts, it is already widely supported in major browsers such as Chrome, Firefox, and Safari, reflecting its rapid adoption for improving web performance beyond traditional WebSockets. This novelty means that there is not much fuzzing work done in this space, presenting an opportunity to contribute to its reliability at an early stage of the lifecycle.

The chosen black-box/gray-box fuzzing approach is advantageous because it treats the target as a black box without requiring access to source code or instrumentation, making it language-agnostic and suitable for testing diverse implementations such as those in Python or Rust, while incorporating gray-box elements through knowledge of the protocol specification to guide mutations. This method is well suited for WebTransport because it allows focusing on message structures and sequences without deep integration into specific runtimes, enabling broader applicability across server libraries. We chose to create a universal fuzzer rather than a language-specific one with deep instrumentation because this supports reuse, avoids dependencies on internal states that may differ across implementations, and aligns with the goal of testing echo servers through external logging and crash detection, ensuring that the fuzzer can evolve with the development of the protocol.

The goals of this thesis include developing a fuzzer using BooFuzz~\cite{boofuzz_docs, boofuzz_github} with message structure mutation and sequence duplication, testing against echo servers, and analyzing results to identify vulnerabilities, addressing challenges such as handling QUIC encryption and the evolving nature of the protocol.

In Chapter~\ref{chap:related} we describe related work and the state of the art in the area of protocol fuzzing. In Chapter~\ref{chap:methodology} we present an overview of the WebTransport protocol, our fuzzing methodology, and the implementation details. In Chapter~\ref{chap:results} we present the results, analysis, and discussion.
